\section{Метод решения}
Программа реализует параллельную версию алгоритма Монте-Карло для вычисления площади окружности с использованием потоков (pthreads) в операционной системе Linux.

{\bfseries Алгоритм работы:}
\begin{enumerate}
    \item Программа принимает параметры командной строки:
    \begin{itemize}
        \item Радиус окружности (-r)
        \item Общее количество точек (-p)
        \item Количество потоков (-t)
        \item Флаг подробного вывода (-v)
    \end{itemize}
    \item Определяется количество ядер процессора для проверки разумности запрошенного количества потоков
    \item Общее количество точек равномерно распределяется между потоками
    \item Создаются и запускаются рабочие потоки, каждый из которых:
    \begin{itemize}
        \item Генерирует случайные точки в квадрате, описывающем окружность
        \item Подсчитывает количество точек, попавших внутрь окружности
        \item Использует независимый генератор случайных чисел (rand\_r) для обеспечения потокобезопасности
    \end{itemize}
    \item Основной поток ожидает завершения всех рабочих потоков
    \item Результаты суммируются с использованием мьютекса для синхронизации доступа к общим переменным
    \item Вычисляется площадь окружности и погрешность по сравнению с теоретическим значением
    \item Выводятся результаты, включая скорость обработки и оценку числа $\pi$
\end{enumerate}

{\bfseries Синхронизация:}
Для предотвращения состояния гонки при обновлении общих переменных используется мьютекс (pthread\_mutex\_t).

{\bfseries Оптимизация:}
\begin{itemize}
    \item Каждый поток использует локальный счётчик для накопления результатов
    \item Синхронизация происходит только при добавлении локальных результатов в общую переменную
    \item Используется потокобезопасный генератор случайных чисел rand\_r
    \item Параллельная генерация случайных точек уменьшает время выполнения
\end{itemize}

\section{Описание программы}
Программа разделена на несколько модулей для обеспечения модульности и повторного использования кода.

{\bfseries Файловая структура:}
\begin{itemize}
    \item \texttt{monte\_carlo.h} — заголовочный файл с объявлениями функций и структур данных
    \item \texttt{monte\_carlo.c} — основная программа с реализацией алгоритма
    \item \texttt{CMakeLists.txt} — файл сборки CMake
\end{itemize}

{\bfseries Основные типы данных:}
\begin{verbatim}
typedef struct {
    long long points_per_thread;  // Количество точек для обработки
    double radius;                // Радиус окружности
    long long points_inside;      // Локальный счётчик точек внутри
    int thread_id;                // Идентификатор потока
    unsigned int seed;            // Сид для генератора случайных чисел
} ThreadData;
\end{verbatim}

{\bfseries Глобальные переменные для синхронизации:}
\begin{verbatim}
pthread_mutex_t mutex;            // Мьютекс для синхронизации
long long total_points_inside;    // Общее количество точек внутри
long long total_points;           // Общее количество обработанных точек
double calculated_area;           // Вычисленная площадь
int active_threads;               // Количество активных потоков
\end{verbatim}

{\bfseries Основные функции:}
\begin{itemize}
    \item \texttt{monte\_carlo\_thread()} — функция потока для генерации и проверки точек
    \item \texttt{parse\_arguments()} — парсинг аргументов командной строки
    \item \texttt{init\_threads()} — создание и инициализация потоков
    \item \texttt{wait\_for\_threads()} — ожидание завершения потоков
    \item \texttt{calculate\_and\_print\_results()} — вычисление и вывод результатов
\end{itemize}

{\bfseries Используемые системные вызовы Linux:}
\begin{itemize}
    \item \texttt{pthread\_create()} — создание потока
    \item \texttt{pthread\_join()} — ожидание завершения потока
    \item \texttt{pthread\_mutex\_init()/pthread\_mutex\_destroy()} — инициализация/уничтожение мьютекса
    \item \texttt{pthread\_mutex\_lock()/pthread\_mutex\_unlock()} — блокировка/разблокировка мьютекса
    \item \texttt{clock\_gettime()} — измерение времени выполнения
    \item \texttt{get\_nprocs()} — получение количества ядер процессора
\end{itemize}

{\bfseries Обработка ошибок:}
Программа проверяет возвращаемые значения всех функций и выводит понятные сообщения об ошибках при:
\begin{itemize}
    \item Некорректных аргументах командной строки
    \item Ошибках создания потоков
    \item Ошибках инициализации мьютекса
\end{itemize}

{\bfseries Демонстрация количества потоков:}
Для демонстрации количества потоков, используемых программой, можно использовать команды:
\begin{itemize}
    \item \texttt{top -H -p <PID>} — отображение потоков процесса
    \item \texttt{ps -eLf | grep monte\_carlo} — список потоков
    \item \texttt{htop} с включённым отображением потоков
\end{itemize}